\section{The big picture: why PDPs aren't enough, and what CDPs do instead}
\label{sec:bigpicture}

\emph{Source: handout \S0; CDP paper (arXiv:2303.04209) \S1 intro. Status:
skeleton from a first read --- reconcile with the paper's own \S1--2 (read
carefully, per the mentor's prescribed order) once read in full. New to
DAGs / the do-operator / backdoor adjustment / potential outcomes? Start
with \S\ref{sec:background} --- everything below assumes it.}

\subsection{The problem}

A fitted model $f$ (any learner) exists, and a picture of how its output
depends on one feature $A$ is wanted, when the features are
\emph{causally related} to each other (not independent).

\textbf{Running example} (used throughout --- reuse it in notes and
questions so everything lines up): age $W \to$ exercise $A \to$ weight
$M \to$ blood pressure $R$; $f$ predicts 10-year cardiovascular risk.

\begin{figure}[h]
\centering
\resizebox{0.92\textwidth}{!}{%
\begin{tikzpicture}[every node/.style={transform shape}]
  \node[defaultnode] (w) {$W$ age};
  \node[defaultnode, right=16mm of w] (a) {$A$ exercise};
  \node[defaultnode, right=16mm of a] (m) {$M$ weight};
  \node[defaultnode, right=16mm of m] (r) {$R$ blood pressure};
  \draw[normaledge] (w) -- (a);
  \draw[normaledge] (a) -- (m);
  \draw[normaledge] (m) -- (r);
  \draw[normaledge] (a) to[bend left=45] (r);
  \draw[normaledge] (w) to[bend right=20] (m);
  \draw[normaledge] (w) to[bend right=38] (r);
\end{tikzpicture}}
\caption{The running example's full causal graph (ECM), reproduced from
handout \S0: every earlier feature is drawn as a direct cause of every later
one (age $\to$ exercise, weight, blood pressure; exercise $\to$ weight,
blood pressure; weight $\to$ blood pressure).}
\label{fig:ecm-base}
\end{figure}

\subsection{The PDP, and why it isn't causally valid}

\textbf{PDP} \reusetag{Friedman 2001}: for each $a$, set everyone's exercise
to $a$, \emph{keep other features as recorded}, average predictions.
Problem: it silently holds weight and blood pressure fixed --- not what
would actually happen if someone exercised more. It answers ``what if only
exercise changed and nothing downstream reacted,'' rarely the causal
question actually wanted.

\medskip
\noindent\textbf{Formally} (Molnar's IML book, ch.~19, following
\reusetag{Friedman 2001}): split all the features the model uses into the
one(s) being plotted, $x_S$, and everything else, $X_C$ ($S$ and $C$
partition every feature $f$ takes as input). The partial dependence
function is
\begin{equation}
\hat f_S(x_S) \;=\; \mathbb{E}_{X_C}\big[\hat f(x_S, X_C)\big]
\;=\; \int \hat f(x_S, X_C)\, dP(X_C).
\label{eq:pdp-def}
\end{equation}
\noindent\textit{In words:} freeze the feature(s) you care about at the
value $x_S$ you want to plot on the $x$-axis, then average the model's
prediction over every possible value the \emph{other} features $X_C$ could
take, weighted by how often that combination actually occurs
($dP(X_C)$ is the real-world distribution of $X_C$). That gives one number
for that $x_S$; repeating it across a range of $x_S$ traces out the curve.

\medskip
\noindent\textbf{Why an integral, not a sum? --- the math one level down.}
$\mathbb{E}_{X_C}[\cdot]$ always means ``probability-weighted average,'' but
what that average looks like as a formula depends on whether $X_C$ is
discrete or continuous.

If $X_C$ could only take finitely many values $c_1,\dots,c_K$ (say, weight
rounded to the nearest kilogram), the weighted average would be an
ordinary sum:
$$\mathbb{E}_{X_C}[\hat f(x_S,X_C)] \;=\; \sum_{k=1}^K \hat f(x_S,c_k)\,
P(X_C=c_k).$$
But weight and blood pressure are \emph{continuous} --- they can take any
real value on a range, so there are infinitely many possible values, and
the probability of hitting any \emph{one exact} value is $0$ (e.g.\
$P(\text{weight}=71.3428\ldots\text{kg exactly})=0$; only \emph{intervals}
of values have positive probability). ``Multiply by $P(X_C=c_k)$ and sum''
breaks down here --- every term would be $0 \times (\text{something})$.

The fix is the same one ordinary calculus uses to define area under a
curve. Probability for a continuous variable is described by a
\textbf{density} $p(x_C)$, where $P(a\le X_C\le b)=\int_a^b p(x_C)\,dx_C$
(a density is a \emph{rate} --- probability per unit of $x_C$ --- not a
probability itself, which is why it's allowed to exceed $1$). Chop the
range of $X_C$ into thin slices of width $\Delta$; on each slice $X_C$ is
approximately discrete, landing in that slice with probability
$\approx p(x_C)\Delta$, so the weighted average over all the slices is
$$\sum_k \hat f(x_S,x_C^{(k)})\, p(x_C^{(k)})\, \Delta
\;\xrightarrow[\ \Delta\to 0\ ]{}\;
\int \hat f(x_S,x_C)\, p(x_C)\, dx_C.$$
That limit --- summing over infinitely many, infinitely thin slices --- is
exactly what an integral sign \emph{means}. It's the same idea as ``area
under a curve is the limit of a sum of thin rectangles'' from a first
calculus course, just with $\hat f(x_S,x_C)\,p(x_C)$ playing the role of
the curve's height.

\emph{Why $dP(X_C)$ instead of $p(x_C)\,dx_C$?} Writing $dP(X_C)$ is a
deliberate generalization: it means ``weighted by however probability is
actually spread out,'' without committing to whether $X_C$ is discrete,
continuous, or a mix. When $X_C$ is continuous, $dP(X_C) = p(x_C)\,dx_C$;
when it's discrete, $\int(\cdot)\,dP(X_C)$ collapses back into the ordinary
sum $\sum_k(\cdot)\,P(X_C{=}c_k)$ above --- one symbol covers both cases,
which is why $P(dw)$ in the notation table (\S\ref{sec:notation}) got the
same ``integrate or sum, whichever applies'' description.

\textbf{Where did $p(x_C^{(k)})\,\Delta$ go in the Monte Carlo sum below?}
Eq.~\eqref{eq:pdp-mc} is $\frac1n\sum_i \hat f(x_S,x_C^{(i)})$ --- no density,
no $\Delta$, every row weighted equally by $1/n$. That's not a different
formula from the slice-sum above; the density weight is still there, just
hidden inside \emph{which points you're summing over}. Here's the
derivation, one step at a time.

\begin{enumerate}
  \item Go back to the fixed grid of slices from before, spaced $\Delta$
    apart, with representative points $x_C^{(1)},x_C^{(2)},\dots$ --- these
    are evenly-spaced, chosen \emph{by you}, regardless of where real data
    lives. That's why each one needed an explicit
    $p(x_C^{(k)})\,\Delta$ multiplier: a grid point in a crowded region of
    the population and a grid point in a near-empty region get exactly one
    term each on the same grid, so probability has to be bolted on by hand
    to tell them apart.
  \item Now overlay your actual dataset on that same grid. Let $N_k$ be the
    number of your $n$ rows whose $x_C^{(i)}$ lands in slice $k$. Group the
    raw (unweighted) sum by which slice each row falls into --- $\hat f$
    barely changes across one thin slice, so every row in slice $k$
    contributes almost the same value, $\hat f(x_S,x_C^{(k)})$:
    $$\sum_{i=1}^n \hat f(x_S,x_C^{(i)}) \;\approx\;
    \sum_k N_k\,\hat f(x_S,x_C^{(k)}).$$
  \item Divide both sides by $n$:
    $$\frac1n\sum_{i=1}^n \hat f(x_S,x_C^{(i)}) \;\approx\;
    \sum_k \underbrace{\frac{N_k}{n}}_{\text{share of your data in slice }k}
    \hat f(x_S,x_C^{(k)}).$$
  \item This is the step that answers the question. $N_k/n$ is the
    \emph{empirical} fraction of your dataset that happens to land in
    slice $k$. If your dataset rows are a representative sample of the
    population (each one drawn from the true distribution of $X_C$), the
    Law of Large Numbers says that fraction converges, as $n$ grows, to the
    slice's \emph{true} probability --- which is exactly what a density
    means:
    $$\frac{N_k}{n} \;\xrightarrow[n\to\infty]{}\; P(X_C\in\text{slice }k)
    \;\approx\; p(x_C^{(k)})\,\Delta.$$
\end{enumerate}
So $p(x_C^{(k)})\,\Delta$ didn't disappear --- it's sitting inside $N_k/n$.
A crowded region of the population (high $p$) automatically contributes
\emph{many} rows to your dataset, so it automatically gets many terms in
the plain, unweighted sum; a rare region contributes few rows and few
terms. The multiplying-by-density step and the counting-more-rows-there
step are doing the identical job --- multiplying explicitly \emph{and}
sampling proportionally to density would double-count the same weighting.
Concretely: 10 people in your data with weight $\approx 70$kg contribute
10 near-identical terms to the raw average, automatically giving that
weight bracket $10\times$ the say of a bracket only 1 person has ---
that \emph{is} the density weighting, just paid for in row-count instead
of in an explicit multiplier. (This relies on the dataset actually being
representative of the population you want the curve to describe --- the
same ``which reference population'' issue flagged in \S\ref{sec:scope}.)

Putting the two limits together --- $\Delta\to 0$ (arbitrarily fine
slices) and $n\to\infty$ (arbitrarily large data) --- turns the slice-sum
into the exact integral, which is why Eq.~\eqref{eq:pdp-mc} is called an
\emph{estimator of} Eq.~\eqref{eq:pdp-def} rather than a separate
definition: same target, approached by letting real data stand in for the
unknown true distribution.

You can't literally integrate over the true, unknown distribution of
$X_C$, so in practice $\hat f_S$ is estimated by \textbf{Monte Carlo
averaging} over the dataset you actually have:
\begin{equation}
\hat f_S(x_S) \;=\; \frac{1}{n}\sum_{i=1}^n \hat f\big(x_S,\, x_C^{(i)}\big).
\label{eq:pdp-mc}
\end{equation}
\noindent\textit{In words:} go through every row $i$ in your dataset, one
at a time. For each row, swap in the fixed value $x_S$ for the feature(s)
you're plotting, but keep that \emph{same row's own, real, observed}
values for every other feature ($x_C^{(i)}$ --- the ``$(i)$'' means ``as
recorded for individual $i$,'' not a made-up value). Run the model once
per row under this swap, then average all $n$ of the resulting
predictions. ``Monte Carlo'' here just means: instead of doing the exact
integral in Eq.~\eqref{eq:pdp-def} (impossible, since you don't know the
true distribution $P(X_C)$), approximate it by an average over samples you
actually have --- your dataset itself stands in for that unknown
distribution, and the approximation gets better the more rows you average
over.

This is exactly the same computation as averaging every individual's ICE
(individual conditional expectation) curve at that $x_S$ --- a PDP curve
\emph{is}, by construction, the pointwise mean of all the ICE curves in
the dataset. And it's precisely the operation already described above in
words (``set everyone's exercise to $a$, keep other features as
recorded, average predictions''): Eq.~\eqref{eq:pdp-mc} is that sentence
written as a formula, with $S=\{A\}$ (exercise) and $C=\{M,R\}$ (weight,
blood pressure) in the running example.

\medskip
\noindent\textbf{A fully worked example, with numbers.} Suppose the
dataset has just $n=5$ people (a toy size, purely so every row can be
written out by hand), and a toy model
$$\hat f(A,M,R) \;=\; 100 - 2A + 0.3M + 0.4R$$
predicting a cardiovascular risk score from (exercise hrs/week $A$,
weight kg $M$, blood pressure $R$) --- a made-up linear formula, invented
only so the arithmetic below is checkable by hand; a real $\hat f$ would
be some fitted ML model with no such closed form.

\begin{table}[h]
\centering
\begin{tabular}{@{}cccc@{}}
\toprule
$i$ & $A_i$ (real) & $M_i$ (real) & $R_i$ (real) \\
\midrule
1 & 2 & 70 & 120 \\
2 & 5 & 65 & 110 \\
3 & 0 & 90 & 140 \\
4 & 3 & 75 & 125 \\
5 & 1 & 85 & 135 \\
\bottomrule
\end{tabular}
\caption{Toy dataset, $n=5$ people, each with their own recorded
$(A_i,M_i,R_i)$.}
\label{tab:pdp-toy}
\end{table}

To compute $\hat f_S(5)$ --- ``what would average predicted risk be if
\emph{everyone} exercised 5 hrs/week'' --- Eq.~\eqref{eq:pdp-mc} says: take
each row in Table~\ref{tab:pdp-toy}, replace $A_i$ with the fixed value
$5$, keep that same row's own $M_i,R_i$ completely untouched, run
$\hat f$, then average all 5 results:

\begin{table}[h]
\centering
\begin{tabular}{@{}ccl@{}}
\toprule
$i$ & $(x_S,\,x_C^{(i)}) = (5,\ M_i,\ R_i)$ & $\hat f(5,M_i,R_i)$ \\
\midrule
1 & $(5,\ 70,\ 120)$ & $100-10+21+48=159.0$ \\
2 & $(5,\ 65,\ 110)$ & $100-10+19.5+44=153.5$ \\
3 & $(5,\ 90,\ 140)$ & $100-10+27+56=173.0$ \\
4 & $(5,\ 75,\ 125)$ & $100-10+22.5+50=162.5$ \\
5 & $(5,\ 85,\ 135)$ & $100-10+25.5+54=169.5$ \\
\bottomrule
\end{tabular}
\caption{The 5 terms actually being averaged in Eq.~\eqref{eq:pdp-mc} at
$x_S=5$. Column 1 (exercise) is the same fixed $5$ in every row --- the
``freeze $x_S$'' part. Weight and blood pressure differ row by row ---
each person's own, untouched $x_C^{(i)}$.}
\label{tab:pdp-toy-eval}
\end{table}

$$\hat f_S(5) \;=\; \frac{159.0+153.5+173.0+162.5+169.5}{5}
\;=\; \frac{817.5}{5} \;=\; 163.5.$$

\noindent\textit{In words:} we are averaging over \textbf{5 numbers, one
per person in the dataset} --- not over 5 possible values of exercise, and
not over some abstract population. Person 3, who in reality never
exercises at all ($A_3=0$, Table~\ref{tab:pdp-toy}), still contributes a
term to this average: their \emph{predicted risk in a pretend world where
they exercised 5 hrs/week but kept their actual weight and blood
pressure} (173.0 --- the single highest of the five, driven entirely by
their real $M_3=90,R_3=140$, since $A$ is fixed at $5$ for every row
here, exercise contributes nothing to the differences between rows).
Repeating the same 5 rows with $x_S=0$ instead of $5$ gives
$\hat f_S(0)=173.5$ instead --- a different final number, because only
column 1 changed while the 5 real $(M_i,R_i)$ pairs stayed exactly the
same. That's the entire mechanism behind \emph{one point} of the PDP
curve; repeating this 5-row average once per value of $a$ traces out the
whole curve.

\textbf{The independence assumption, made concrete.} Eq.~\eqref{eq:pdp-mc}
pairs the fixed $x_S=a$ with \emph{every} row's own $x_C^{(i)}$, including
rows whose real exercise level was nothing like $a$. That's fine only if
exercise and (weight, blood pressure) are unrelated --- but they aren't
(that's the whole running example). Concretely, in
Table~\ref{tab:pdp-toy-eval}: Person 3 really never exercises
($A_3=0$) and really carries a high weight/BP profile ($M_3=90,R_3=140$)
--- plausibly \emph{because} they don't exercise. Eq.~\eqref{eq:pdp-mc}
still uses their $(M_3,R_3)$ as a stand-in for ``someone who exercises 5
hrs/week,'' as if that weight/BP combination were a realistic thing to see
at high exercise, when in this population it may never actually occur
there.
This is the same mechanism as the confounding problem in
\S\ref{sec:background} --- it's why the PDP ``isn't causally valid'' named
in this subsection's title, and it's exactly what the CDP framework below
is built to fix.

\bigskip
\noindent\textbf{Aside: PDP-based feature importance}
\reusetag{Greenwell, Boehmke \& McCarthy 2018}. Not part of the CDP
framework itself, but built directly on the PDP curve just defined, and
worth having on hand since ``does this curve actually vary, or is it
flat?'' is a question that comes back repeatedly later.

\emph{The idea in one sentence:} a flat PDP curve means the feature barely
changes the prediction (not important); a curve that swings up and down a
lot means the feature matters a lot (important). Feature importance turns
``how much does the curve vary'' into a single number.

\textbf{Numerical features.} Evaluate the PDP curve $\hat f_S$ at its $K$
\emph{unique observed values} $x_S^{(1)},\dots,x_S^{(K)}$ (literally every
distinct value the feature actually takes in the dataset, not an
arbitrary grid), then measure how spread out those $K$ curve heights are
around their own average:
\begin{equation}
I(x_S) \;=\; \frac{1}{K-1}\sum_{k=1}^K
\left(\hat f_S\big(x_S^{(k)}\big) - \frac{1}{K}\sum_{k=1}^K
\hat f_S\big(x_S^{(k)}\big)\right)^{\!2}.
\label{eq:pdp-importance-numeric}
\end{equation}
\noindent\textit{In words:} the inner term
$\frac1K\sum_k \hat f_S(x_S^{(k)})$ is just the \emph{average height} of
the curve across its $K$ points --- the flat reference line you'd see if
the feature had zero effect. For each point, subtract that average and
square the result: a point on a flat curve sits close to the average, so
this is close to $0$; a point far from the average (the curve is
swinging there) gives a large squared number. Add up all $K$ squared
deviations and divide by $K-1$. This is nothing more exotic than the
ordinary \textbf{sample variance formula}
($s^2=\frac1{n-1}\sum(y_i-\bar y)^2$, same $K-1$ denominator for the same
reason) applied to the $K$ curve heights instead of $K$ raw data points
--- ``importance'' here literally \emph{means} ``variance of the PDP
curve.''

\emph{Worked example, continuing the toy dataset above.} Since
$\hat f_S(a) = 100-2a+73.5 = 173.5-2a$ in the running toy (the average of
$0.3M_i+0.4R_i$ over the 5 people works out to $73.5$), the $K=5$ unique
real exercise values $\{0,1,2,3,5\}$ give curve heights
$\{173.5,\,171.5,\,169.5,\,167.5,\,163.5\}$, averaging to $169.1$. Squared
deviations: $4.4^2,2.4^2,0.4^2,(-1.6)^2,(-5.6)^2 =
19.36,\,5.76,\,0.16,\,2.56,\,31.36$, summing to $59.2$. So
$I(A) = 59.2/4 = 14.8$. If exercise instead had \emph{no} effect on
$\hat f$ at all, all 5 curve heights would be identical, every squared
deviation would be $0$, and $I(A)$ would come out exactly $0$ --- that's
the whole point of the measure.

\textbf{Categorical features.} Categories don't sit on a number line, so
``deviation from the mean'' isn't quite the natural notion; a cruder,
more robust measure is used instead --- the \textbf{range rule}:
\begin{equation}
I(X_S) \;=\; \frac{\max_k\big(\hat f_S(x_S^{(k)})\big)
- \min_k\big(\hat f_S(x_S^{(k)})\big)}{4}.
\label{eq:pdp-importance-categorical}
\end{equation}
\noindent\textit{In words:} take the PDP's highest value across the $K$
categories, subtract its lowest value, divide by $4$. \emph{Why divide by
$4$?} A rule of thumb borrowed from the normal distribution: about $95\%$
of a normal distribution's mass sits within $\pm 2$ standard deviations of
its mean --- a band $4$ standard deviations wide. So if all you know is a
variable's range (max$-$min) and want a rough guess at its standard
deviation, range$/4$ is a reasonable back-of-envelope estimate
($\text{range}\approx 4\sigma \Rightarrow \sigma\approx
\text{range}/4$). It's deliberately crude and typically
\emph{underestimates} the true spread --- but with only a handful of
categories ($K$ small), there's little basis for a real variance estimate
anyway, which is exactly why the cruder range-based version replaces
Eq.~\eqref{eq:pdp-importance-numeric} here.

\emph{Worked example (toy, illustrative numbers, following Molnar's own
season/bike-rental example):} PDP-predicted bike rentals by season,
Spring$=150$, Summer$=180$, Fall$=165$, Winter$=90$. Range
$=180-90=90$, so $I(\text{season}) = 90/4 = 22.5$.

\textbf{Two things this measure is \emph{not}.} (1) It only sees the
feature's \emph{main effect} --- the shape of its own PDP curve --- and
is blind to \emph{interactions}: a feature that matters only by
interacting with another feature (e.g.\ flipping sign depending on a
second feature's value, the heterogeneous-effects problem from the
recall questions above) can have a nearly flat, low-$I$ PDP curve while
still being genuinely important to the model --- permutation feature
importance would catch that; this measure would miss it. (2) It weights
every \emph{unique value} equally, regardless of how many rows share it
--- a value only one person in the dataset has counts exactly as much as
a value shared by hundreds, unlike the PDP curve itself
(Eq.~\eqref{eq:pdp-mc}), which already weights by how common a value's
neighborhood is via the Monte Carlo average.

\subsection{The fix: use the ECM to define ``changing \texorpdfstring{$A$}{A}''}

\textbf{ECM} = ``explanatory causal model'' --- the analyst's DAG over the
\emph{features}. It is an assumption brought to the explanation, not a claim
verified by the model. Given an ECM, there is more than one way to ask ``how
does output depend on $A$'' --- there are (at least) four, depending on
whether downstream features are allowed to respond:

\begin{table}[h]
\centering
\begin{tabular}{@{}p{2.6cm}p{5.2cm}p{5.6cm}@{}}
\toprule
\textbf{Plot} & \textbf{Question} & \textbf{In the running example} \\
\midrule
TDP (total) & set $A=a$; let everything downstream respond &
  weight and BP move \\
NDDP (natural direct) & set $A=a$; hold downstream features at their
  \emph{natural} values & \textbf{this is exactly the PDP} (handout's
  Thm.\ 2.10) \\
NIDP (natural indirect) & keep your own $A$; move downstream features to
  where they'd be under $A=a$ & your exercise, but the weight/BP you'd have
  at exercise $a$ \\
PCDP (partially controlled) & set $A=a$ \emph{and} fix another feature &
  exercise $=a$, weight $=m$ \\
\bottomrule
\end{tabular}
\caption{The four named dependence plots.}
\label{tab:four-plots}
\end{table}

The genuinely new fact worth internalizing early: $\mathbf{NDDP = PDP}$. The
plot everyone already uses is a special case of this framework --- the
``direct effect only'' member of the family, not a separate method.

\subsection{Constructed outcome --- the trick that makes this identifiable}

Treat the model's output as a variable: $Y^* = f(X)$. Then the
\textbf{average TDP is a plain back-door adjustment}:
\begin{equation}
  \theta(a) \;=\; E[f(X) \mid \operatorname{do}(A=a)]
  \;=\; E_W\big[\, E[Y^* \mid A=a, W] \,\big], \qquad W = \text{parents of } A.
  \label{eq:tdp-backdoor}
\end{equation}
\emph{In words:} the average model output you'd get from forcing everyone
to exercise level $a$ ($\theta(a)$, left-hand side) equals a two-step
average on the right: first, within each age group $w$, average the
model's prediction over the people who actually exercise at level $a$
(inner $E[Y^*\mid A=a,W]$); then average those age-group numbers together
using the real population mix of ages ($E_W[\cdot]$) --- the same
``average over the marginal, not the conditional'' move as
Eq.~\ref{eq:backdoor}, now applied to the model's output instead of a raw
outcome.
\reusetag{Pearl 2009, Thm.\ 3.2.2; Zhao \& Hastie 2021 (theorem)}; \newtag~is
what this implies specifically for CDPs. The equations for weight and blood
pressure are never needed: $\operatorname{do}(A=a)$ leaves the parents'
distribution alone and leaves every downstream conditional given $(a,w)$
alone. This is \emph{why} \S\ref{sec:identification} can get away with only
needing the parents plus a ``reduced form,'' not the full fitted ECM.

\subsection{The picture that organizes the whole family}

\reusetag{Shpitser \& Tchetgen Tchetgen 2016} Add the model's output $\hat Y$
to the graph as a child of every feature. Then TDP/PCDP are ordinary node
interventions ($\operatorname{do}(A=a)$, $\operatorname{do}(A=a,C=c)$), while
NDDP/NIDP are \textbf{edge} interventions (set $a$ on the arrow into $\hat
Y$ only, vs.\ on the arrows into the child features only). Normally, one
node ``acting two ways at once'' needs cross-world assumptions data cannot
check --- but here $\hat Y = f(X)$ is a \emph{known function}, so nothing
cross-world is actually invoked. That is the structural reason NDDP reduces
to an ordinary average (the PDP) and NIDP reduces to an ordinary
single-world quantity (\S\ref{sec:identification}).

\begin{figure}[htbp]
\centering
\begin{subfigure}[t]{0.47\textwidth}
\centering
\begin{tikzpicture}[scale=0.62, every node/.style={transform shape}]
  \node[defaultnode] (w) at (0,1.6) {$W$};
  \node[accentnode] (a) at (2.4,1.6) {$A{=}a$};
  \node[respondnode] (m) at (4.8,1.6) {$M$};
  \node[respondnode] (r) at (7.2,1.6) {$R$};
  \node[targetnode] (y) at (3.6,-0.8) {$\hat Y$};
  \draw[severededge] (w) -- (a);
  \draw[normaledge] (a) -- (m);
  \draw[normaledge] (m) -- (r);
  \draw[normaledge] (a) -- (y);
  \draw[normaledge] (m) -- (y);
  \draw[normaledge] (r) -- (y);
\end{tikzpicture}
\caption{TDP --- total: \texttt{do(A=a)} cuts the arrow into $A$; $M$ and $R$
respond through the normal mechanism; everything reaches $\hat Y$.}
\label{fig:tdp}
\end{subfigure}
\hfill
\begin{subfigure}[t]{0.47\textwidth}
\centering
\begin{tikzpicture}[scale=0.62, every node/.style={transform shape}]
  \node[defaultnode] (w) at (0,1.6) {$W$};
  \node[accentnode] (a) at (2.4,1.6) {$A$};
  \node[naturalnode] (m) at (4.8,1.6) {$M$};
  \node[naturalnode] (r) at (7.2,1.6) {$R$};
  \node[targetnode] (y) at (3.6,-0.8) {$\hat Y$};
  \draw[normaledge] (w) -- (a);
  \draw[severededge] (a) -- (m) node[midway, above, font=\tiny, text=black!55] {obs.\ $A$};
  \draw[normaledge] (m) -- (r);
  \draw[accentedge] (a) -- (y) node[midway, left, font=\tiny, text=accentc] {$a$};
  \draw[normaledge] (m) -- (y);
  \draw[normaledge] (r) -- (y);
\end{tikzpicture}
\caption{NDDP $=$ PDP --- natural direct: only the arrow into $\hat Y$
carries $a$; $M$ and $R$ are left at their natural (recorded) values.}
\label{fig:nddp}
\end{subfigure}

\medskip

\begin{subfigure}[t]{0.47\textwidth}
\centering
\begin{tikzpicture}[scale=0.62, every node/.style={transform shape}]
  \node[defaultnode] (w) at (0,1.6) {$W$};
  \node[accentnode] (a) at (2.4,1.6) {$A$};
  \node[respondnode] (m) at (4.8,1.6) {$M$};
  \node[respondnode] (r) at (7.2,1.6) {$R$};
  \node[targetnode] (y) at (3.6,-0.8) {$\hat Y$};
  \draw[normaledge] (w) -- (a);
  \draw[accentedge] (a) -- (m) node[midway, above, font=\tiny, text=accentc] {$a$};
  \draw[normaledge] (m) -- (r);
  \draw[severededge] (a) -- (y) node[midway, left, font=\tiny, text=black!55] {obs.\ $A$};
  \draw[normaledge] (m) -- (y);
  \draw[normaledge] (r) -- (y);
\end{tikzpicture}
\caption{NIDP --- natural indirect: the arrow into $M$ carries $a$ (so $M,R$
move to their counterfactual values); the direct arrow into $\hat Y$ keeps
the observed $A$. Identified as Fulcher et al.'s ``bridge mean.''}
\label{fig:nidp}
\end{subfigure}
\hfill
\begin{subfigure}[t]{0.47\textwidth}
\centering
\begin{tikzpicture}[scale=0.62, every node/.style={transform shape}]
  \node[defaultnode] (w) at (0,1.6) {$W$};
  \node[accentnode] (a) at (2.4,1.6) {$A{=}a$};
  \node[accentnode] (m) at (4.8,1.6) {$M{=}m$};
  \node[respondnode] (r) at (7.2,1.6) {$R$};
  \node[targetnode] (y) at (3.6,-0.8) {$\hat Y$};
  \draw[severededge] (w) -- (a);
  \draw[severededge] (a) -- (m);
  \draw[normaledge] (m) -- (r);
  \draw[normaledge] (a) -- (y);
  \draw[normaledge] (m) -- (y);
  \draw[normaledge] (r) -- (y);
\end{tikzpicture}
\caption{PCDP --- partially controlled: \texttt{do(A=a,M=m)} cuts both
incoming arrows; $R$ still responds to the fixed $M$; both $a$ and $m$ reach
$\hat Y$.}
\label{fig:pcdp}
\end{subfigure}
\caption{The four dependence plots as different interventions on the same
explanatory causal model. \emph{Simplified for legibility}: the full ECM
(Figure~\ref{fig:ecm-base}, handout \S0) also has direct edges $W\to M$,
$W \to R$, $A \to R$ among the features, and --- since $\hat Y$ is a child
of every feature --- a direct $W \to \hat Y$ edge; all four are omitted here
so the direct/indirect mechanism stays readable. Blue $=$ set to the
counterfactual value; amber $=$ responds/recomputed; gray $=$ kept at its
natural (recorded) value; dashed $=$ edge severed, or carrying the natural
value instead of $a$.}
\label{fig:four-plots}
\end{figure}

\noindent\textbf{Legend:}
\begin{itemize}[leftmargin=1.6em, itemsep=2pt]
  \item[\textcolor{accentc}{$\bullet$}] set to the counterfactual value
  \item[\textcolor{respondc}{$\bullet$}] responds / recomputed
  \item[\textcolor{naturalc}{$\bullet$}] kept at natural (recorded) value
  \item[-----] active edge (solid)
  \item[- - -] severed edge, or an edge carrying the natural value instead
    of $a$
\end{itemize}

\begin{quote}
\textbf{Direct + indirect $\neq$ total.} \newtag, modest.
\[
  \theta(a) - B \;=\; [\mathrm{NDDP}(a) - B] + [\mathrm{NIDP}(a) - B] + J(a),
\]
\emph{In words:} take how far the total curve is from the overall average
prediction ($\theta(a) - B$, left side). You might expect that gap to split
cleanly into ``how much of it is the direct effect'' plus ``how much of it
is the indirect effect'' (the two bracketed terms on the right). It doesn't,
in general --- there's a leftover piece $J(a)$ that has to be added to make
the two sides match, and $J(a)$ is exactly zero only in the special case
where $f$ treats $A$ and the downstream features additively (no
interaction between them).

Here $B$ is the average prediction and $J$ is an interaction remainder that
vanishes only when $f$ is additive in $A$ and the downstream features. Toy
counterexample: $A \sim \mathrm{Bern}(p)$, $R(a) = a$, $f = a \cdot r$; the
total contrast is $1$, but the direct and indirect contrasts are each only
$p$. \textbf{Don't assume direct $+$ indirect sums to total} --- $J$ should
be plotted alongside them, not ignored.
\end{quote}

\begin{openquestions}
  \item Work through why exactly the parents' distribution is untouched by
    $\operatorname{do}(A=a)$ --- is this just the definition of ``parents,''
    or does it need the ECM to have no unmeasured confounding among
    features?
  \item Read arXiv \S1--2 in full (the mentor's prescribed order has this
    \emph{before} rereading the handout) for the paper's own formal
    definitions, then compare notation to the handout's.
\end{openquestions}
